\section{The large relative gaps}
\label{app:positive-boundary}

The electronic proof appendix \textit{Positive-exponent Boundary}
contains the contour estimates and the exact finite checks used in
this section. The following argument specifies the ranges in which
each of those proofs is needed.

\begin{lemma}
\label{lem:positive-boundary}
The pair \(P_{N-1,a,d},P_{N,a,d}\) is coprime for every \(N\ge2\)
on \(a=1,d\ge3\), and on \(a\ge2,d\ge3a\).
\end{lemma}

\begin{proof}
Proposition~\ref{prop:small-terminals} covers \(N\le24\).
Thus we may impose the ratio bound of Lemma~\ref{lem:ratio-range}
throughout this proof.

For \(a=1\), write
\[
 Q_{m,\gamma,d}(x)=\sum_{k=0}^m\binom{m}{k}
                        (\gamma)_{dk}x^{m-k}.
\]
At terminal \(N\) the actual shapes are \(N\) and \(N+1\).
Coefficient comparison gives, with \(\gamma=N\),
\[
 Q_{N,\gamma,d}'=NQ_{N-1,\gamma,d},\qquad
 \gamma Q_{N,\gamma+1,d}
  =(\gamma+dN)Q_{N,\gamma,d}-dxQ_{N,\gamma,d}'.
\]
Consequently the actual pair has a common root exactly when
\(Q_{N,N,d}\) has a multiple root. The prime-power proof in the
electronic appendix gives squarefreeness at
\(d=3,4,5,7,8,9\) for this integral shape. The separate
gap-six contour proof gives it at \(d=6\) for \(N\ge10000\).

For \(d\ge10\), the lower-one contour proof gives squarefreeness of
\(Q_{N,N,d}\) whenever \(N\ge\max(5000,128(d+1))\).
Its two contour contributions have separated logarithmic
derivatives at a zero, so the derivative of their sum is nonzero.
For \(a\ge2,d\ge3a\), the weighted contour proof in the electronic
appendix gives coprimality for \(N\ge\max(4000,128(a+d))\).
Here one must compare two integrals with the actual moment orders.
With \(w=A/(BN^d)\) and \(g(x)=x^a(w+x^d)\), their integrands
after the common dilation are \(g(x)^N e^{-Nx}\) and
\(g(x)^{N-1}e^{-Nx}\). Write \(I_U,I_L\) and \(J_U,J_L\)
for their two contour contributions. At a zero of the first
integral, \(I_U=-I_L\ne0\). The verified packet estimates at the
two selected saddles \(x_U,x_L\) give
\[
 \left|\frac{J_U}{I_U}\right|
 <\frac{24\sqrt{N/(N-1)}}{|g(x_U)|},\qquad
 \left|\frac{J_L}{I_L}\right|
 >\frac{\sqrt{N/(N-1)}}{5|g(x_L)|},\qquad
 \left|\frac{g(x_U)}{g(x_L)}\right|>150 .
\]
Thus \(|J_U|<|J_L|\) and the second integral is nonzero.
The proof estimates the extra amplitude \(e^{-x}\) in the
order-\(N-1\) contour integral when the phase is written as
\(\log g(x)-x\).

If \(N<\max(4000,128b)\), where \(b=a+d\le4d/3\), a
potential failure has \(d\le24\). Indeed \(s\ge3d/4\).
For \(N\ge200\) and \(d\ge25\), one has
\[
 N<512d/3,\qquad \frac{3d}{4}>2\log(512d/3)+2 .
\]
The second inequality holds at \(d=25\) and its margin increases.
For \(N<200\) the ratio bound is smaller still.

For \(a=1,d\ge10\) below the contour threshold,
a potential failure has \(d\le19\).
For \(20\le d\le38\) the threshold is \(5000\), and
\[
 \frac{d^2}{d+1}\ge\frac{400}{21}>2\log(5000)+2 .
\]
For \(d\ge39\) the threshold is \(128(d+1)\), and the
same comparison holds at \(39\) with an increasing margin.
The orders \(N<200\) again satisfy the weaker ratio bound
with room.

The supplemental finite checker therefore needs only the families
\[
\begin{array}{c|c|c}
a&d&\text{terminals }N\\ \hline
1&6&25\le N<10000\\
1&10\le d\le19&25\le N<5000\\
2\le a\le\lfloor d/3\rfloor&6\le d\le24&
                25\le N<\max(4000,128(a+d)).
\end{array}
\]
Before any further exclusions this is \(338\,071\) cases.
Every case satisfies \(a<d(N-2)\), and a possible common factor
has \(s<21\) by Lemma~\ref{lem:ratio-range}. The block-prime
support and Newton allocation arguments in
Lemmas~\ref{lem:prime-support} and~\ref{lem:newton-alloc}
apply here with these hypotheses. The stronger degree bound
is used only when its starting-point condition is verified
for that degree. In the subsequent Newton check, each odd
prime \(p>aN\) satisfies \(p>N,p>d\) and \(p^2>(a+d)N\)
throughout these finite families. Thus the sparse floor
ordinates used in Section~\ref{sec:arithmetic} are exact
also in this range.

The supplemental program \texttt{verify\_positive\_boundary.cpp}
retains \(326\,146\) cases after imposing the uniform bound
\(d^2<21(a+d)\). It tries the
first available sparse Newton test, applies the prime-block
bounds at the remaining degrees, and then uses additional
primes as needed. The exact exclusions are
\[
\begin{array}{l|r|r|r}
\text{Family}&\text{Supports}&\text{Prime blocks}&\text{Newton polygons}\\ \hline
a\ge2,\ d\ge3a &266\,421&12\,901&253\,520\\
a=1,\ 10\le d\le19&49\,750&966&48\,784\\
a=1,\ d=6&9\,975&535&9\,440
\end{array}
\]
Every possible even degree is excluded, and no modular fallback
is needed.
\end{proof}
